\documentclass[UTF8,a4paper,10pt]{ctexart}
\usepackage[a4paper,top=14mm,bottom=16mm,left=16mm,right=16mm,footskip=8mm]{geometry}
\usepackage{amsmath,amssymb,booktabs,tabularx,array,enumitem,fancyhdr,lastpage,xcolor,hyperref,multirow}
\setCJKmainfont{SimSun}\setCJKsansfont{SimHei}\setmainfont{Times New Roman}
\hypersetup{hidelinks,unicode=true}\pagestyle{fancy}\fancyhf{}\fancyfoot[C]{第\thepage 页\quad 共\pageref*{LastPage}页}\renewcommand{\headrulewidth}{0pt}
\renewcommand{\arraystretch}{1.2}\setlist[enumerate]{leftmargin=2.2em,itemsep=.24em,topsep=.2em}
\newcommand{\sect}[1]{\par\smallskip\noindent\fcolorbox{black}{black!4}{\parbox{\dimexpr\linewidth-2\fboxsep-2\fboxrule}{\sffamily\bfseries #1}}\par\smallskip}
\newcommand{\opts}[1]{\par\hspace*{.5em}\begin{tabularx}{\dimexpr\linewidth-.5em}{@{}XXXX@{}}#1\end{tabularx}\par}
\begin{document}\small
\begin{center}{\zihao{2}\sffamily\bfseries 湖南工商大学课程考核试卷【问卷】}\par{\zihao{3}\bfseries 统计学B（A卷）}\end{center}
\noindent\begin{tabularx}{\linewidth}{@{}lXlX@{}}学分：&3&考核学期：&2025--2026学年度第一学期\\考核形式：&闭卷&年级、专业：&2024级各专业全校基础课\\时量：&120分钟&&\end{tabularx}
\noindent 考生务必将答案写在答卷上，问卷上答题一律无效；考生可用普通计算器，结果保留4位小数。
\sect{一、单选题（每小题1分，共10分）}
\begin{enumerate}
\item 通过有限的种子发芽实验来估计整批种子的发芽率。这种方法属于（）。\opts{A.推断统计学&B.描述统计学&C.数学&D.逻辑学}
\item 一个统计总体（）。\opts{A.只能有一个标志&B.可以有多个标志&C.只能有一个指标&D.可以有多个指标}
\item 下列调查方式中属于全面调查的是（）。\opts{A.普查&B.重点调查&C.典型调查&D.抽样调查}
\item 下面属于变量数列的是（）。\opts{A.职工按性别分组&B.电站按发电量分组&C.企业按所有制类型分组&D.企业按经济部门分组}
\item 计算年距增长量或年距增长率的目的是为了消除下列哪一因素变动的影响（）。\opts{A.长期趋势&B.季节变动&C.循环变动&D.随机波动}
\item 反映总体分布集中趋势的是（）。\opts{A.总量指标&B.相对指标&C.平均指标&D.标志变异指标}
\item 累计增长量等于相应时期各个逐期增长量的（）。\opts{A.和&B.差&C.积&D.商}
\item 现象在其发展变化过程中所呈现出来的那种持续上升或持续下降的趋势是（）。\opts{A.长期趋势&B.季节变动&C.循环变动&D.不规则变动}
\item 拉氏指数所采用的同度量因素是固定在（）。\opts{A.基期&B.报告期&C.假定期&D.任意时期}
\item 下列抽样方式中，抽样误差最小的是（其他条件相同）（）。\opts{A.纯随机抽样&B.等距抽样&C.分层抽样&D.整群抽样}
\end{enumerate}
\sect{二、判断题（每小题1分，共10分）}
\begin{enumerate}[start=11]
\item 标志说明总体单位的特征，而指标则说明总体的特征。（）
\item 质量指标不能用数值表示，数量指标能用数值表示。（）
\item 一般来说，离散型变量组距数列，相邻组的组限必须重叠。（）
\item 用组中值代表各组的一般水平的假定条件是各组组距相等。（）
\item 不管是按品质标志分组，还是按数量标志分组，都能形成变量数列。（）
\item 产品合格率属于结构相对数。（）
\item 如果两个数列的全距相同，那么它们的离散程度也相同。（）
\item 如果总体分布没有明显的集中趋势，众数也可以不存在。（）
\item 统计量是总体的数量特征，参数是样本的数量特征。（）
\item 回归系数可用来判断现象之间的相关方向。（）
\end{enumerate}
\sect{三、填空题（每小题4分，每空2分，共20分）}
\begin{enumerate}[start=21]
\item 某企业劳动生产率计划比上年提高20\%，实际提高30\%。则该企业劳动生产率的计划完成程度为\underline{\hspace{2.5cm}}，是否完成了任务（选填是或否）\underline{\hspace{1.5cm}}。
\item 适度偏态下某数列的众数$M_o=85$，中位数$M_e=79$。则该数列的算术平均数的大约值为\underline{\hspace{2cm}}，该数列的分布特征为（选填左偏或右偏）\underline{\hspace{1.5cm}}。
\item 已知$n=10,\sum x=351,\sum y=400,\sum x^2=13915,\sum y^2=18350,\sum xy=15832$，则$x,y$的相关系数$r$为\underline{\hspace{2cm}}，相关程度为（选填低度、显著或高度相关）\underline{\hspace{2cm}}。
\item 某数列的算术平均数为16.3，各变量值平方的平均数为272.9。则该数列方差为\underline{\hspace{2cm}}，标准差系数为\underline{\hspace{2cm}}。
\item 某企业2010--2015年间利润年均发展速度为115\%，2022年利润较2015年翻了2番，2024年利润是2022年的1.5倍。则2010--2024年间该企业利润年均发展速度为\underline{\hspace{2cm}}，年均增长率为\underline{\hspace{2cm}}。
\end{enumerate}
\sect{四、计算题}
\begin{enumerate}[start=26]
\item 某企业生产一批产品2000件，采取不重复抽样方法随机抽选300件进行检测，发现其中合格品285件，若置信概率为95\%（$t=1.96$）。试计算：
\begin{enumerate}[label=（\arabic*）,leftmargin=2.8em]\item 该批产品合格率的抽样平均误差；（4分）\item 该批产品合格率的抽样极限误差；（4分）\item 该批产品合格率的置信区间。（4分）\end{enumerate}
\item 某企业产品产量情况如下表（单位：吨），请完成下列计算。
\begin{center}\begin{tabular}{crrrrr}\toprule 年度&第一季度&第二季度&第三季度&第四季度&总平均或合计\\\midrule2022年&275&495&190&130&1090\\2023年&320&450&200&90&1060\\2024年&290&510&170&140&1110\\平均&&&&&\\季节指数&&&&&\\2025年&300&&&&\\\bottomrule\end{tabular}\end{center}
\begin{enumerate}[label=（\arabic*）,leftmargin=2.8em]\item 计算各季度平均值、三年的季度总平均值，填入表格；（4分）\item 计算季节指数、季节指数之和，填入表格；（4分）\item 2025年第一季度若为300吨，计算该年各季及全年产量，填入表格。（4分）\end{enumerate}
\item 商场A、B、C三种商品2023--2024年的销售资料如下表（单位：元），请完成下列计算。
\begin{center}\begin{tabular}{crrrrrrr}\toprule\multirow{2}{*}{品种}&\multicolumn{2}{c}{价格}&\multicolumn{2}{c}{销量}&\multicolumn{3}{c}{销售额}\\\cmidrule(lr){2-3}\cmidrule(lr){4-5}\cmidrule(lr){6-8}&$p_0$&$p_1$&$q_0$&$q_1$&$p_0q_0$&$p_1q_1$&$p_0q_1$\\\midrule A&160&130&1000&1200&&&\\B&70&50&1500&1600&&&\\C&60&90&5200&4900&&&\\合计&&&&&&&\\\bottomrule\end{tabular}\end{center}
\begin{enumerate}[label=（\arabic*）,leftmargin=2.8em]\item 计算相应的销售额数据，填入表格；（3分）\item 计算三种商品的价格总指数、销量总指数、销售额总指数；（6分）\item 计算价格和销量变动对销售额变动的影响额。（3分）\end{enumerate}
\end{enumerate}
\sect{五、综合题（每小题12分，共24分）}
\begin{enumerate}[start=29]
\item 某企业50名工人产量数据如下表（单位：件）。
\begin{center}\begin{tabular}{*{10}{r}}38&45&49&52&58&60&62&62&63&63\\65&65&67&69&70&71&71&73&73&75\\75&75&76&76&77&78&78&78&79&79\\79&79&81&83&85&86&86&87&87&88\\88&89&89&90&91&91&93&93&96&99\end{tabular}\end{center}
\begin{enumerate}[label=（\arabic*）,leftmargin=2.8em]\item 根据原始数据编制组距式变量数列，填入如下表格：（4分）
\begin{center}\begin{tabular}{ccccc}\toprule 产量&组中值$x$&次数$f$&频率&$xf$\\\midrule60以下&&&&\\60--70&&&&\\70--80&&&&\\80--90&&&&\\90以上&&&&\\合计&&&&\\\bottomrule\end{tabular}\end{center}
\item 根据变量数列计算产量的加权算术平均数、中位数和众数；（6分）\item 根据算术平均数、中位数和众数分析该数列的分布特征。（2分）\end{enumerate}
\item 我国居民2005--2024年人均国内生产总值$x$与城镇居民人均可支配收入$y$数据（单位：百元）的Excel回归分析结果如下表。
\begin{center}\begin{tabular}{lr}\multicolumn{2}{c}{回归统计}\\\toprule Multiple R&0.9977\\R Square&0.9953\\Adjusted R Square&0.9951\\标准误差&4.4188\\观测值&20\\\bottomrule\end{tabular}\par\smallskip
\begin{tabular}{lrrrrr}\multicolumn{6}{c}{方差分析}\\\toprule &df&SS&MS&F&Significance F\\回归分析&1&74720.48&74720.48&3826.74&$2.00835\mathrm E{-22}$\\残差&18&351.47&B&&\\总计&A&75071.94&&&\\\bottomrule\end{tabular}\par\smallskip
\begin{tabular}{lrrrr}\toprule &Coefficients&标准误差&t Stat&P-value\\Intercept&$-9.7093$&2.2647&$-4.2873$&0.000443448\\X Variable 1&0.2424&0.0039&61.8607&$2.00835\mathrm E{-22}$\\\bottomrule\end{tabular}\end{center}
\begin{enumerate}[label=（\arabic*）,leftmargin=2.8em]\item 写出相关系数和估计标准误差；（2分）\item 写出表格中的阴影部分（A、B）的数值；（4分）\item 写出人均国内生产总值与城镇居民人均可支配收入的一元线性回归方程；（4分）\item 若2025年人均国内生产总值1000百元，预测该年城镇居民人均可支配收入。（2分）\end{enumerate}
\end{enumerate}
\newpage\begin{center}{\zihao{2}\sffamily\bfseries 参考答案}\end{center}
\sect{客观题与填空题}
选择题：1--5 ADABB；6--10 CAAAC。判断题：对、错、错、错、错、对、错、对、错、对。\par
21. 108.33\%，是；22. 73，左偏；23. 0.926，高度相关；24. 7.21，16.48\%；25. 114.87\%，14.87\%。
\sect{大题校对说明}
第26--30题已按原题恢复全部题干、分问、原始数据表和Excel回归表。源文件的答案页仅给出客观题及第21、22、23、25题答案；其余大题没有可可靠转录的完整官方解答，故未擅自补写。
\end{document}
